\chapter*{Notation reference}
\addcontentsline{toc}{chapter}{Notation reference}
\setcounter{section}{0}\label{sec:root:notation-reference:1}

A lookup table connecting the mathematical notation used in this book's prose to the corresponding Lean syntax used in its code. This page is a quick reference, not something to read start to finish --- every row is built up and explained properly the first time it is needed, at the chapter cited in the last column.

Two related pages cover notation this one deliberately leaves out: the \hyperref[sec:root:tactic-and-library-reference:1]{tactic and library reference} (tactics and Mathlib declaration names) and the \hyperref[sec:root:lambda-calculus-dictionary:1]{λ-calculus / type theory to Lean dictionary} (formal type-theory notation from the optional ``Mathematical reading'' boxes). This page is just the ordinary logic/algebra symbols used throughout the main text and code.

\section{Logic and quantifiers}\label{sec:root:notation-reference:2}

{\small\begin{tabular}{
  >{\raggedright\arraybackslash}p{(\linewidth - 8\tabcolsep) * \real{0.3400}}
  >{\raggedright\arraybackslash}p{(\linewidth - 8\tabcolsep) * \real{0.2000}}
  >{\raggedright\arraybackslash}p{(\linewidth - 8\tabcolsep) * \real{0.2800}}
  >{\raggedright\arraybackslash}p{(\linewidth - 8\tabcolsep) * \real{0.1800}}
}
\toprule
Meaning & Math notation & Lean syntax & First appears \\
\midrule
Function type / implication & \(A \to B\) & \texttt{A → B} & Chapter 1 \\
Universal quantifier (``for all'') & \(\forall x, P\, x\) & \texttt{∀ x, P x} & Chapter 3 \\
Existential quantifier (``there exists'') & \(\exists x, P\, x\) & \texttt{∃ x, P x} & Chapter 3 \\
Unique existence (``there exists a unique'') & \(\exists!\, x, P\, x\) & no single token --- witnessed by supplying the value and a proof it is the only one & Chapter 1 §4 \\
Set/type membership & \(x \in A\) & \texttt{x ∈ A} & Chapter 1 \\
Negation & \(\neg P\) & \texttt{¬P} & Chapter 3 \\
Conjunction (``and'') & \(P \wedge Q\) & \texttt{P ∧ Q} (\texttt{And}) & Chapter 3 \\
Disjunction (``or'') & \(P \vee Q\) & \texttt{P ∨ Q} (\texttt{Or}) & Chapter 3 \\
Not equal & \(a \neq b\) & \texttt{a ≠ b} & Chapter 3 \\
Turnstile (``the goal to prove'') & \(\Gamma \vdash P\) & the goal-state display (not typed by the user) & Chapter 4 \\
\bottomrule
\end{tabular}}


\section{Algebra, structure, and diagrams}\label{sec:root:notation-reference:3}

{\small\begin{tabular}{
  >{\raggedright\arraybackslash}p{(\linewidth - 8\tabcolsep) * \real{0.3400}}
  >{\raggedright\arraybackslash}p{(\linewidth - 8\tabcolsep) * \real{0.2000}}
  >{\raggedright\arraybackslash}p{(\linewidth - 8\tabcolsep) * \real{0.2800}}
  >{\raggedright\arraybackslash}p{(\linewidth - 8\tabcolsep) * \real{0.1800}}
}
\toprule
Meaning & Math notation & Lean syntax & First appears \\
\midrule
Definitional equality & \(t \equiv t'\) & \texttt{rfl} closes the goal & Chapter 5 §4 \\
Isomorphism / equivalence & \(A \simeq B\) & \texttt{A ≃ B} (\texttt{Equiv}) & Chapter 10 \\
Anonymous-constructor pairing & \(\langle a, b \rangle\) & \texttt{⟨a, b⟩} & Chapter 2 §1 \\
Function composition & \(g \circ f\) & \texttt{g ∘ f} & Chapter 1 \\
Scalar/group action, or a generic infix operation & \(a \cdot b\) & \texttt{a • b} (\texttt{SMul}) & Chapter 10 \\
Inverse & \(a^{-1}\) & \texttt{a⁻¹} & Chapter 6 \\
Lambda abstraction (``sends to'') & \(x \mapsto e\) & \texttt{fun x => e} & Chapter 1 \\
Divisibility & \(a \mid b\) & \texttt{a ∣ b} & Chapter 9 \\
Subset & \(A \subseteq B\) & \texttt{A ⊆ B} & Chapter 10 \\
Cartesian product & \(A \times B\) & \texttt{A × B} & Chapter 1 \\
Long/derivation arrow (diagrams) & \(A \longrightarrow B\) & \texttt{⟶} (diagram labels only, not ordinary code) & Chapter 1 §4 \\
Projections out of a product & \(\pi_X, \pi_Y\) & \texttt{.1}/\texttt{.2}, or \texttt{.fst}/\texttt{.snd} & Chapter 1 §4 \\
\bottomrule
\end{tabular}}

